\documentclass{article}
\usepackage[utf8]{inputenc}
\usepackage{amsmath,amssymb,amsthm}
\usepackage{booktabs}
\usepackage{graphicx}
\usepackage{hyperref}
\usepackage{natbib}

\newtheorem{theorem}{Theorem}
\newtheorem{definition}{Definition}
\newtheorem{observation}{Observation}

\title{Why Your Knowledge Graph Uncertainty Is Lying to You:\\The Coverage Blind Spot}

\author{Anonymous}

\begin{document}

\maketitle

\begin{abstract}
Knowledge graph embedding (KGE) models are increasingly deployed in high-stakes applications, yet their uncertainty estimates fail silently on a critical class of queries: \emph{novel contexts}---familiar entities appearing in unseen relational contexts. We demonstrate that \textbf{widely-used uncertainty methods} (energy scores, probabilistic embeddings, ensembles, MC dropout) achieve near-random performance (AUROC 0.46--0.62) on novel context detection across standard benchmarks. Strikingly, Energy is \emph{most confident} on zero-coverage queries---60\% of its top-100 confident predictions have no training evidence for the queried context. While these predictions often happen to be correct via compositional generalization, the model is operating without direct evidence. We analyze why this failure is structural: these methods compute entity-level uncertainty that cannot distinguish relational contexts. The only solution is explicit coverage tracking---a simple hash table lookup. Critically, we show that coverage-based detection and prediction-quality estimation are \emph{fundamentally different tasks}: coverage flags queries without evidence, while confidence predicts answer correctness. Practitioners need both. We provide concrete guidelines: (1) track entity-relation coverage, (2) flag zero-coverage queries for review, (3) use coverage for OOD detection and confidence for quality estimation.
\end{abstract}

\section{Introduction}

Knowledge graph embedding models power critical applications from drug discovery~\citep{bonner2022review} to financial fraud detection. When these models make predictions, practitioners need to know: \emph{should I trust this answer?}

The standard approach is uncertainty quantification---estimating how confident the model is. Methods range from energy-based scores to probabilistic embeddings to deep ensembles. A common assumption is that higher uncertainty indicates less reliable predictions.

\textbf{We show this assumption fails on a common query type.}

Consider querying a KG about Albert Einstein: ``What instrument did Einstein play?'' A model trained on physics-related facts about Einstein has never seen him in a musical context. Yet existing uncertainty methods report \emph{high confidence}---the model ``knows'' Einstein well (high entity frequency) so it must know the answer.

This is the \textbf{coverage blind spot}: uncertainty methods that ignore relational context cannot detect when a familiar entity appears in an unfamiliar role.

\subsection{Contributions}

\begin{enumerate}
    \item \textbf{Empirical finding}: We show widely-used KG uncertainty methods fail on novel context detection (Section~\ref{sec:experiments}). Energy is most confident on zero-coverage queries---60\% of its top-100 have no training evidence.

    \item \textbf{Analysis}: We explain why this failure occurs: existing methods compute entity-level uncertainty that is constant across relational contexts (Section~\ref{sec:analysis}).

    \item \textbf{Key distinction}: Coverage-based uncertainty and confidence-based quality estimation serve \emph{different purposes}. Coverage asks ``have I seen this context?'' while confidence asks ``is this answer correct?'' Practitioners need both (Section~\ref{sec:distinction}).

    \item \textbf{Practical guidelines}: Concrete recommendations for deploying KG systems with proper uncertainty handling (Section~\ref{sec:guidelines}).
\end{enumerate}

\section{The Coverage Blind Spot}

\subsection{Problem Setup}

Let $\mathcal{G} = (\mathcal{E}, \mathcal{R}, \mathcal{T})$ be a knowledge graph with entities $\mathcal{E}$, relations $\mathcal{R}$, and triples $\mathcal{T} \subseteq \mathcal{E} \times \mathcal{R} \times \mathcal{E}$.

\begin{definition}[Coverage]
The coverage indicator $\text{cov}(e, r) = 1$ if entity $e$ appears with relation $r$ in training, else 0.
\end{definition}

\begin{definition}[Novel Context]
A test triple $(h, r, t)$ is a \emph{novel context} if $\text{cov}(h, r) = 0$ or $\text{cov}(t, r) = 0$---at least one entity appears with a relation it was never trained on.
\end{definition}

Novel contexts are common: 32\% of FB15k-237 test triples, 16\% of YAGO3-10, 31\% of ICEWS14.

\subsection{Why Existing Methods Fail}
\label{sec:analysis}

We examine four widely-used uncertainty methods:
\begin{itemize}
    \item \textbf{Energy}: $U = -f(h, r, t)$ (negative score)
    \item \textbf{UKGE}: $U = \sigma^2(h) + \sigma^2(t)$ (learned entity variances)
    \item \textbf{Deep Ensemble}: $U = \text{Var}[f_1, \ldots, f_K]$ (model disagreement)
    \item \textbf{MC Dropout}: $U = \text{Var}[\text{samples}]$ (dropout sampling)
\end{itemize}

\begin{observation}[Architectural Limitation]
All four methods compute uncertainty as a function that does not condition on the specific (entity, relation) pair. UKGE assigns $\sigma^2(e)$ per entity regardless of relation. Energy, Ensemble, and MC Dropout measure confidence in the \emph{answer} but not whether the \emph{question} is in-distribution.
\end{observation}

This architectural choice means the same entity receives the same base uncertainty whether queried about familiar or unfamiliar relations.

\section{Experiments}
\label{sec:experiments}

\subsection{Setup}

We evaluate on four benchmarks: FB15k-237, WN18RR, YAGO3-10, and ICEWS14.

\textbf{Task}: Detect novel context triples (binary classification, measured by AUROC).

\textbf{Note on Table~\ref{tab:main}}: Energy achieves higher AUROC on WN18RR (0.78) and ICEWS14 (0.89). These datasets have structural properties (WN18RR: only 11 relations; ICEWS14: temporal patterns) where Energy's confidence happens to correlate with coverage. This correlation is dataset-specific, not a property of the method. On FB15k-237 and YAGO3-10, Energy achieves 0.62 and 0.46 respectively.

\subsection{Main Results}

\begin{table}[h]
\centering
\caption{Novel context detection AUROC. Learned methods show inconsistent, often near-random performance. Coverage tracking achieves perfect detection by construction.}
\label{tab:main}
\begin{tabular}{lcccc}
\toprule
Method & FB15k-237 & WN18RR & YAGO3-10 & ICEWS14 \\
\midrule
Energy & 0.62 & 0.78$^\dagger$ & 0.46 & 0.89$^\dagger$ \\
UKGE & 0.50 & 0.50 & 0.50 & 0.50 \\
Deep Ensemble & 0.47 & 0.48 & -- & -- \\
MC Dropout & 0.53 & 0.52 & -- & -- \\
\midrule
Coverage-Only & 1.00 & 1.00 & 1.00 & 1.00 \\
\bottomrule
\end{tabular}
\end{table}

$^\dagger$Higher AUROC due to dataset-specific correlation, not method design (see text).

Key findings:
\begin{itemize}
    \item \textbf{UKGE achieves exactly 0.50} (random) across all datasets---entity-level variance provides zero signal for novel context detection
    \item \textbf{Energy achieves 0.46 on YAGO3-10}---worse than random, meaning it is \emph{more confident on OOD queries}
    \item \textbf{Coverage tracking achieves 1.00} by construction---it directly answers the question ``is this context covered?''
\end{itemize}

\subsection{The Key Finding: Confident Without Evidence}

We analyze Energy's most confident predictions on FB15k-237.

\begin{table}[h]
\centering
\caption{Energy's most confident predictions on FB15k-237. Top-100 confident predictions are 60\% zero-coverage, yet achieve 95\% Hits@10.}
\label{tab:confident}
\begin{tabular}{lccc}
\toprule
 & Top-100 & Top-1000 & All Test \\
\midrule
Zero-coverage rate & \textbf{60\%} & 47\% & 32\% \\
Hits@10 & 95\% & 94\% & 40\% \\
\bottomrule
\end{tabular}
\end{table}

A surprising finding: \textbf{zero-coverage predictions are often correct} (57.6\% Hits@10 vs 32.5\% for covered). This reinforces our key point---coverage and correctness are \emph{orthogonal}:

\begin{itemize}
    \item Energy is confident on zero-coverage queries (60\% of top-100)
    \item These predictions happen to be correct via compositional generalization
    \item But the model has \textbf{no direct evidence}---it is guessing correctly by chance
\end{itemize}

\textbf{Why this matters}: A correct guess without evidence is not trustworthy. The same model will fail on the next novel context that doesn't follow the pattern. Coverage flags queries where the model is operating outside its training distribution, regardless of whether it guesses correctly.

\section{The Only Solution: Coverage Tracking}
\label{sec:solution}

\subsection{Implementation}

Coverage tracking requires a simple data structure:
\begin{verbatim}
coverage = set()
for (h, r, t) in train_triples:
    coverage.add((h, r))
    coverage.add((t, r))

def is_novel_context(h, r, t):
    return (h, r) not in coverage or
           (t, r) not in coverage
\end{verbatim}

\subsection{Scalability}

\begin{table}[h]
\centering
\caption{Coverage storage requirements. Based on actual coverage density from each dataset.}
\begin{tabular}{lcc}
\toprule
Dataset & Hash Table & Bloom (1\% FPR) \\
\midrule
FB15k-237 (3.4M pairs) & 26 MB & 5 MB \\
YAGO3-10 (4.6M pairs) & 35 MB & 7 MB \\
Wikidata-scale (est.) & 174 GB & 33 GB \\
\bottomrule
\end{tabular}
\end{table}

For most deployments, a hash table suffices. For web-scale KGs, Bloom filters provide 5$\times$ compression with 1\% false positive rate.

\section{Critical Distinction: Detection vs.\ Quality}
\label{sec:distinction}

A natural question: if coverage detects OOD queries, does it also predict answer quality?

\textbf{No. These are fundamentally different tasks.}

\begin{table}[h]
\centering
\caption{Correlation between uncertainty and prediction error (Spearman $\rho$).}
\begin{tabular}{lcc}
\toprule
Method & $\rho$ (error) & Interpretation \\
\midrule
Energy & +0.42 & Predicts error \\
Coverage & $-$0.20 & Anti-correlated \\
\bottomrule
\end{tabular}
\end{table}

\textbf{Why the negative correlation?} Zero-coverage queries are often ``easy'' patterns the model can guess correctly via compositional generalization. Coverage absence does not mean the model will be wrong---it means the model has \emph{no direct evidence}. A correct guess without evidence is not trustworthy.

This reveals two distinct questions:
\begin{itemize}
    \item \textbf{``Have I seen this context?''} $\to$ Coverage tracking
    \item \textbf{``Is this answer correct?''} $\to$ Confidence scores
\end{itemize}

\textbf{Practitioners need both signals.} A zero-coverage query that happens to be answered correctly is still a reliability concern---the model got lucky, and may fail on the next similar query.

\section{Practical Guidelines}
\label{sec:guidelines}

Based on our findings, we recommend:

\begin{enumerate}
    \item \textbf{Always track coverage}: Maintain an (entity, relation) lookup table alongside your KG model.

    \item \textbf{Flag zero-coverage queries}: When $\text{cov}(h,r) = 0$ or $\text{cov}(t,r) = 0$, the model has no direct evidence. Flag for human review or indicate lower reliability.

    \item \textbf{Use both signals}: Coverage for OOD detection, confidence for quality estimation. Don't conflate them.

    \item \textbf{Report stratified metrics}: Evaluate separately on covered and novel-context splits. Aggregate metrics hide the coverage blind spot.
\end{enumerate}

\section{Related Work}

\textbf{KG uncertainty quantification.} UKGE~\citep{chen2019embedding} learns entity-level confidence. BEUrRE~\citep{chen2021probabilistic} uses box embeddings. Both compute entity-level uncertainty that cannot distinguish relational contexts.

\textbf{OOD detection.} Energy-based OOD detection~\citep{liu2020energy} works well in vision where the OOD signal is in the input distribution. In KGs, the OOD signal is in the \emph{relational context}, which existing architectures ignore.

\textbf{Relation-aware methods.} R-GCN and CompGCN use relation-specific message passing for \emph{scoring}, but not for \emph{uncertainty}. Extending these to relation-conditioned uncertainty is an interesting direction, though coverage tracking provides a simpler solution.

\section{Conclusion}

We identified a critical blind spot in KG uncertainty quantification: widely-used methods fail to detect novel contexts, achieving near-random AUROC. Energy is most confident on zero-coverage queries---operating without direct evidence even when predictions happen to be correct.

The solution is simple: track (entity, relation) coverage. A hash table lookup succeeds where sophisticated neural methods fail.

Critically, we showed that coverage-based OOD detection and confidence-based quality estimation are \emph{different tasks}. Coverage flags queries without direct evidence; confidence predicts answer correctness. Both matter for reliable deployment.

Our findings suggest practitioners should:
\begin{itemize}
    \item Deploy coverage tracking alongside any KG model
    \item Use coverage for ``should I trust this?'' and confidence for ``how accurate is this?''
    \item Report metrics stratified by coverage status
\end{itemize}

The gap between ``model knows this entity'' and ``model knows this entity \emph{in this context}'' is the coverage blind spot. Closing it requires no new algorithms---just awareness and a lookup table.

\bibliographystyle{plainnat}
\bibliography{references}

\end{document}
